\subsection{Empirical anchors in single-photon timing measurements}
\label{sec:empiricalAnchors}

The operational quantities above are not tied to a particular detector technology, but classical SPAD measurements provide concrete examples of the stochastic timing structure that motivates them. In early high-resolution time-correlated photon-counting work, Cova \emph{et al.} described the measured record explicitly as a histogram of photon-detection delays, with separate contributions from source synchronization, optical-pulse width, detector timing jitter, and associated timing electronics \cite{Cova1989}. They also noted that convolution analysis can estimate decay constants substantially shorter than the instrumental FWHM, illustrating why a single response width is not itself an information metric.

Measured detector IRFs are likewise not generally characterized by one width. Spinelli \emph{et al.} reported SPAD responses with a regular Gaussian-like fast component associated with absorption in the high-field depletion region and slower diffusion tails associated with carriers generated in neutral material \cite{Spinelli1998}. They emphasized that tailing and secondary bumps can obscure weak or delayed temporal structure even when the FWHM is competitive. Their double-junction SPAD reduced the residual tail to a small fast exponential component, whereas comparison detectors could exhibit substantially longer exponential diffusion tails. This is precisely the type of shape dependence retained by $G(\omega)$ and $B_{\rm FI}$ but discarded by a scalar FWHM.

Spatial and readout variables also provide physical examples of potentially informative marks or of downstream coarse graining. Lacaita and Mastrapasqua showed that, in shallow-junction SPADs, the avalanche-current leading edge and discriminator-crossing delay depend on where the avalanche is initiated and on detector diameter, because the multiplication region must spread across the active area \cite{LacaitaMastrapasqua1990}. Later measurements by Assanelli \emph{et al.} found that timing jitter in re-engineered SPADs depends on both injection position and discriminator threshold, and that fixing the injection position does not remove all timing dispersion because avalanche-propagation statistics remain \cite{Assanelli2011}. These observations motivate the distinction in Sec.~\ref{sec:markGradient} between retaining a latency-predictive primary-event variable and discarding it in a thresholded or otherwise coarse-grained record.

Finally, reach-through APDs provide a direct example of an unresolved fast stochastic internal process. Lacaita \emph{et al.} found that photon-assisted lateral avalanche spreading, driven by secondary photons emitted by hot carriers and reabsorbed elsewhere in the device, accounts quantitatively for the observed avalanche transient and a large fraction of the measured timing jitter \cite{Lacaita1993}. The example is technology-specific, but the lesson is general: a microscopic stochastic transition mechanism can set the conditional timing law even when its internal trajectory is not directly observed. The empirical examples in this subsection motivate the resource variables used here; they are not assumptions required by the theorem stack.
